\section{LLM annotation protocol}
\label{supp:s1-llm}

This note documents the LLM-based annotation pipeline used to label Weibo posts for emotional arousal and risk content.

\subsection{Model and infrastructure}

\begin{itemize}
  \item Model: \texttt{Qwen/Qwen3-8B}.
  \item Serving: vLLM via an OpenAI-compatible Chat Completions API.
  \item Client: \texttt{openai} Python client, enforcing a single JSON object output using \texttt{response\_format=\{"type":"json\_object"\}}.
\end{itemize}

\subsection{Inference settings}

\begin{itemize}
  \item Temperature: 0.1
  \item Max tokens: 512
\end{itemize}

\subsection{Annotation tasks}

Each post is labeled along two dimensions:
\begin{itemize}
  \item Emotional arousal (\texttt{emotion\_class}): \texttt{H} (high), \texttt{M} (medium/neutral), \texttt{L} (low).
  \item Risk content (\texttt{risk\_class}): \texttt{risk} vs.\ \texttt{norisk}.
\end{itemize}

\subsection{Prompt template and output schema}

We use a structured prompt that defines label criteria and requires the model to output a single JSON object (no additional text). The implementation is in \texttt{src/empirical/llm\_annotator.py} and is invoked by \texttt{scripts/run\_new\_annotation.py}.

\subsubsection{System prompt (English translation for readability)}

\begin{verbatim}
You are an expert in social media content analysis. Your task is to label each Weibo post
along two dimensions: emotional arousal and risk content. Output exactly one JSON object
and nothing else (no chain-of-thought, no markdown).

Emotion arousal (emotion_class):
- H (high arousal): strong emotions (anger, fear/panic, sarcasm, insults, highly emotional accusations).
- M (medium/neutral): factual or calm discussion, neutral reposts/comments, supportive messages, explanations.
- L (low arousal): negative but not intense emotions (anxiety, worry, confusion, helplessness, sadness).

Risk content (risk_class):
Core rule: does the post convey that COVID-19/long COVID has negative health impacts?
- risk: describes symptoms/health changes, emphasizes severity/long-term/irreversibility, warns about risks,
        reports cases/studies of sequelae, questions "no sequelae" claims, or raises risk-focused questions.
- norisk: emphasizes recovery/controllability, reassuring official explanations, criticizes fear-mongering,
          success stories stressing recovery, or content unrelated to long-COVID health risks.

Edge cases:
- Very short/hashtag-only posts with no substance: default to M + norisk, confidence 0.5.
- Light tone does not imply norisk: joking symptom descriptions are still risk.

Output JSON format (and only this):
{"emotion_class": "H"|"M"|"L",
 "emotion_confidence": 0.0-1.0,
 "risk_class": "risk"|"norisk",
 "risk_confidence": 0.0-1.0,
 "reasoning": "brief justification"}
\end{verbatim}

\subsubsection{User prompt}

\begin{verbatim}
Please analyze the following Weibo post:

---
{text}
---

Return the classification result in JSON format as specified.
\end{verbatim}

\subsubsection{Output fields}

The model must output a single JSON object with the following fields:
\begin{itemize}
  \item \texttt{emotion\_class}: one of \texttt{H}, \texttt{M}, \texttt{L}.
  \item \texttt{emotion\_confidence}: a float in $[0,1]$.
  \item \texttt{risk\_class}: one of \texttt{risk}, \texttt{norisk}.
  \item \texttt{risk\_confidence}: a float in $[0,1]$.
  \item \texttt{reasoning}: a short justification string.
\end{itemize}

\subsection{Validation}

We compared LLM-assigned labels to manual ratings on a manually rated subset ($n=3{,}000$). Overall exact-match agreement (accuracy) between LLM and human labels was 83\%.

